\subsubsection{A single dyadic finite-mode defect test}
\label{subsec:dyadicdefect}
The central-point reduction and the adaptive theta cutoff
can be combined into a single sequence of nonnegative scalar
quantities. This eliminates the need to take separate limits
in the penalty size, Laguerre order and number of theta modes.
It does not by itself establish the required boundedness.

Retain the notation of Theorems~\ref{thm:onepointcriterion}
and~\ref{thm:adaptivepenalty}. In particular, $T=\tfrac12$,
$k_n=(n+1)(2n+1)$ and
\[
\mathcal N_n=
\int_0^T\bigl(P_{T-s}[-U_n(s,\cdot)]_+\bigr)(0)\,ds,
\quad
b_n=1+T(A_n+M_n).
\]
Here $A_n=\sup|U_n|$ and
$M_n=\sup|k_nU_{n+1}|$ on the heat slab,
so $0\le\mathcal N_n\le T A_n$.
For each $C\ge1$, denote by
$\mathcal E_{n,C}^{(M)}(T,0)$
the finite-theta-mode central positive-part penalty
defined in the preceding section.
Let $m_n(C)$ be the explicit cutoff prescribed by
\eqref{eq:jointcutoff}, with its previously proved bound
\[
\big|\mathcal E_{n,C}^{(m_n(C))}(T,0)
-\mathcal E_{n,C}(T,0)\big|<1.
\]

\begin{theorem}[Quantitative dyadic defect limit]
\label{thm:dyadicdefect}
Define the nonnegative infinite-series defect
and its finite adaptive approximations by
\begin{align}
\mathscr D&=\sum_{n=1}^\infty
2^{-n}\frac{\mathcal N_n}{b_n},
\label{eq:defectD}\\
S_k&=\sum_{n=1}^k2^{-n}
\frac{\mathcal E_{n,2^k}^{(m_n(2^k))}(T,0)}{b_n},
\qquad k\ge1.
\label{eq:defectSk}
\end{align}
The series defining $\mathscr D$ converges absolutely,
with $0\le\mathscr D\le1$, and for every $k\ge1$
one has the \emph{explicit unconditional} estimate
\begin{equation}
\boxed{\left|2^{-k}S_k-\mathscr D\right|
<2^{1-k}.}
\label{eq:dyadicrate}
\end{equation}
Moreover, the following conditions are equivalent:
\begin{equation}
\boxed{
\mathrm{RH}\quad\Longleftrightarrow\quad
\mathscr D=0\quad\Longleftrightarrow\quad
S_k<1\ \ \text{for every integer }k\ge1.
}
\label{eq:dyadicRH}
\end{equation}
In particular, a rigorously certified value $S_k\ge1$
at one finite stage would refute RH.
No such value is asserted here.
\end{theorem}
\begin{proof}
By the central penalty sandwich
\eqref{eq:onepointsandwich},
\[
\big|\mathcal E_{n,C}(T,0)-C\mathcal N_n\big|
\le TM_n.
\]
The adaptive estimate then implies
\[
\big|\mathcal E_{n,C}^{(m_n(C))}(T,0)
-C\mathcal N_n\big|<TM_n+1\le b_n.
\]
Set $C=2^k$ and sum over $1\le n\le k$:
\begin{align*}
\left|2^{-k}S_k-
\sum_{n=1}^k2^{-n}\frac{\mathcal N_n}{b_n}\right|
&<2^{-k}\sum_{n=1}^k2^{-n}
<2^{-k}.
\end{align*}
Meanwhile $0\le\mathcal N_n/b_n\le1$, so
the omitted infinite tail is no larger than
$\sum_{n>k}2^{-n}=2^{-k}$.
Combining the bounds gives
\eqref{eq:dyadicrate}.

Strict positivity of the Gaussian heat kernel,
the central-point theorem and the
Laguerre--P\'olya characterization prove that
RH is equivalent to $\mathcal N_n=0$ at every
order. Every summand in $\mathscr D$ is
nonnegative, so this is equivalent to
$\mathscr D=0$.

If RH holds, the true $U_n$ are nonnegative,
hence $\mathcal E_{n,C}(T,0)\le TM_n$.
The adaptive error gives
$\mathcal E_{n,C}^{(m_n(C))}(T,0)
<TM_n+1\le b_n$,
and consequently
$S_k<\sum_{n=1}^k2^{-n}<1$.

Conversely, if RH fails, then $\mathscr D>0$.
Equation~\eqref{eq:dyadicrate} implies
$S_k>2^k\mathscr D-2$.
For all sufficiently large $k$,
this lower bound exceeds $1$.
Thus $S_k<1$ for all $k$ implies RH.
\end{proof}

\paragraph{What this reduction has, and has not, achieved.}
Theorem~\ref{thm:dyadicdefect} is a quantitative
reorganization of the preceding exact heat-penalty
identities. It replaces an uncontrolled interchange
of three limits by one joint dyadic sequence, with
an explicit error $2^{1-k}$.
The cutoff $m_n(2^k)$ depends on $n$ and $k$
but is finite and grows as $O_n(\sqrt{k})$.
Each $S_k$ involves finitely many orders and
finite theta sums, although rigorous evaluation
of its heat integrals over all Fourier heights
still requires separate analytical and interval
control. The difficult global assertion is
\emph{not} convergence of the approximation:
it is $\mathscr D=0$, or equivalently the
unproved bound $S_k<1$ for every $k$.
Neither positivity of the underlying theta
kernel nor the finite-zero disk certificates
establish that inequality.

The rational inequalities in the proof are
independently replayed in
\path{certificates/xi_dyadic_penalty_defect_exact.py}
and the normalization and threshold gates are
formalized in \path{XiKernel/DyadicPenalty.lean}.
These tests do not formalize the infinite
series, the heat integrals or RH end to end;
no new solution of the Riemann hypothesis
is claimed.
